% Fichier généré automatiquement par tools/generate_corrections.py.
% Source : DYN/DYN-05-Methode/09_RT_RSG/09_RT_RSG.tex
% Statut : complet (3/3 question(s) renseignée(s), 0 reprise(s)).
\begin{corrigebox}[Corrigé extrait de la banque]
\CorrigeQuestion{1}
\begin{marginfigure}
\includegraphics[width=\linewidth]{09_RT_RSG_cor_01}
\end{marginfigure}
\CorrigeQuestion{2}
Le système posède deux mobilités : 
\begin{itemize}
\item translation de 1 par rapport à 2 ($\lambda$);
\item rotation de l'ensemble \{1+2\} autour du point $I$ (le roulement sans glissement permet d'écrire une relation entre la rotation de paramètre $\theta$ et le déplacement suivant $\vi{0}$.
\end{itemize}

On en déduit la stratégie suivante : 
\begin{itemize}
\item Première loi de mouvement : 
\begin{itemize}
\item on isole 2,
\item BAME : 
\begin{itemize}
\item $\torseurstat{T}{1}{2}$, 
\item $\torseurstat{T}{1_{\text{ressort}}}{2}$ ($\vectf{1}{2}\cdot \vi{1}= 0$ et $\vectf{1_{\text{ressort}}}{2}\cdot \vi{1}= 0$)
\item $\torseurstat{T}{\text{Pesanteur}}{2}$;
\end{itemize}
\item on réalise un théorème de la résultante dynamique en projection suivant $\vi{1}$. 
\end{itemize}
\item Seconde loi de mouvement : 
\begin{itemize}
\item on isole \{1+2\};
\item BAME :
\begin{itemize} 
\item $\torseurstat{T}{0}{1}$ ($\vectm{I}{0}{1}\cdot \vk{0}= 0$), 
\item $\torseurstat{T}{\text{Pesanteur}}{1}$,
\item  $\torseurstat{T}{\text{Pesanteur}}{2}$.
\end{itemize}
\item on réalise un théorème du moment dynamique en $I$ en projection suivant $\vk{0}$. 
\end{itemize}
\end{itemize}
\CorrigeQuestion{3}
On montre que :
\begin{itemize}
\item $\vectv{B}{2}{0} = \lambdap\vi{1} + \thetap\left(\lambda(t)\vj{1}-R\vi{0} \right)$
\item $\vectg{B}{2}{0}  = \lambdapp(t)\vi{1} 
+ \thetapp(t)\left(\lambda(t)\vj{1}-R\vi{0} \right)
+ \thetap(t)\left(2\lambdap(t)\vj{1}-\lambda(t)\thetap\vi{1} \right)
$
\item $\vectrd{2}{0} =  m_2 \left( \lambdapp(t)\vi{1} +\lambdap(t)\thetap\vj{1} 
+ \thetapp(t)\left(\lambda(t)\vj{1}-R\vi{0} \right)
+ \thetap(t)\left(\lambdap(t)\vj{1}-\lambda(t)\thetap\vi{1} \right)\right)$
\item $\vectmd{I}{2}{0} \cdot \vect{k_0} =     
-m_2 R\lambdapp(t)\cos\theta 
+m_2 R\lambdap(t)\thetap\sin\theta 
 +m_2 R^2\thetapp(t) 
+ m_2 R\thetap(t)\lambdap(t)\sin\theta
+m_2 R\thetap^2(t)\lambda(t)\cos\theta 
+ 2 m_2  \lambda(t)  \lambdap(t)\thetap
+  m_2  \thetapp(t) \lambda(t)^2 +C_2\thetapp$
\end{itemize}
\end{corrigebox}
